\documentclass[12pt]{article}
\usepackage{graphicx} % Required for inserting images
\usepackage{hyperref}
\usepackage{makecell}
\usepackage{makeidx}
\begin{document}

\graphicspath{ {../templates/img/} }

\newcommand{\firstPage}{
    \thispagestyle{empty}
    \begin{figure}
    \centering
    \includegraphics[scale=0.7]{Swellfish_logo.png}
    \end{figure}
    \author{
        \date{}
        \href{mailto://swellfish14@gmail.com}{swellfish14@gmail.com} \\
    } 
}
\section{Introduzione}
Lo scopo del documento è quello di fornire un insieme di regole, strumenti e stadard di qualità necessari per creare un Way of Working condiviso e efficacemente impiegabile da tutti i membri del gruppo.
Questo documento conterrà tutte le comvenzioni da utilizzare per lo sviluppo del progetto e verrà aggiornato in maniera incrementale durante tutte le fasi dello sviluppo, pertanto è da considerarsi come work-in-progress.


\section{Modalità di sviluppo}
Data la necessità di effettuare un lavoro sostanzioso e condiviso, come tutti gli sviluppatori software, il team ha deciso di utilizzare una repository pubblica su GitHub per rendere agevole la collaborazione e la visione delle modifiche apportate alle varie parti del progetto.
E' stata quindi creata un'organizzazione chiamata SWEllFish e al suo interno sono state create delle specifiche repository:

\begin{itemize}
    \item SWELLFish: questa repo conterrà tutto il codice necessario per il progetto e i vari branch necessari per la suddivisione dei compiti
    \item Documentazione: questa repo contiene tutti i documenti utilizzabili sia dagli sviluppatori che dai committenti per avere una visione chiara delle modifiche apportate al progetto.
\end{itemize}

\section{Contatti con il committente/proponente}
Per semplificare la gestione dei contatti e delle eventuali richieste con l'azienda committente e con il Professore Vardenega, è stata creata una mail, \href{swellfish14@gmail.com}{} utilizzabile da tutti i componenti del gruppo.
L'uso di questa casella è stato regolamentato con un insieme di regole condivise dagli sviluppatori, pertanto un messaggio nasce dalla necessità di uno o più componenti del gruppo di avere delucidazioni da parte del committente o per sottoporre una modifica in attesa di accettazione.
Il canale telegram verrà usato per chiedere informazioni in modo informale con il proponente oppure per richiedere un confronto veloce 
I colloqui vengono richiesti via mail oppure via telegram, tramite l'apposito gruppoo creato con il responsabile incaricato dell'azienda proponente.

\section{Sviluppo}
La scadenza per la presentazione delle candidature dei capitolati è stata fissata per il 27/03/2023, pertanto il gruppo prevede di iniziare i lavori di analisi/sviluppo indicativamente il 07/04/2023, lasciando una settimana di "buco" per organizzare al meglio i compiti e le tempistiche.

\section{Organizzazione dei File}
Tutta la documentazione inerente al progetto sarà incusa nella repository "Documentazione" all'interno dell'organizzazione chiamata "SwellFish".
I documenti di interesse per il proponente/committente si trovano in Documentazione/esterni, mentre quelli necessari al team di sviluppo si trovano in Documentazione/interni. I verbali infine sono contenuti in Documentazione/esterni/verbali.

Sono state applicate delle specifiche regole da adottare per la nomenclatura dei verbali e ogni documento dovrà rispettare la seguente struttura: Verbale_data_azienda/nome_docente.

\section{Processi di supporto}
\subsection{Documentazione}
Ogni documento che verrà redatto attraverserà diversi processi, ogniuno dei quali vedrà impegnati uno o più componenti del gruppo.

Fasi:
\begin{itemize}
    \item Pianificazione: il file viene ideato e discusso con i membri del team sulla base delle necessità che deve soddisfare
    \item Stesura: il testo e il contenuto del documento viene realizzato dal redattore
    \item Revisione: il revisore incaricato si occupa di rileggerlo e correggere eventuali errori e/o imprecisioni, controllando che siano state rispettate le convenzioni stabilite in questo documento.
    \item Approvazione: il responsabile designato controlla che il documento sia corretto e coerente.
\end{itemize}
\subsection{Struttura del documento}
Per facilitare la produzione di documentazione e ridurre i tempi necessari per la redazione, è stato ideato un template generico da utilizzare.
Questo documento ha una struttura ben specifica:
\begin{itemize}
    \item Nome del gruppo
    \item Logo 
    \item Componenti del gruppo
    \item Data/Periodo di realizzazione
    \item Tabella con ruoli dei componenti del gruppo e finalità d'uso del documento
    \item Indice
\end{itemize}
\section{Versionamento}
Il VCS scelto dal gruppo per facilitare la gestione delle modifiche è Git, mediante l'utilizzo di GitHub.
A questo scopo è stata creata un'organizzazione denominata "SwellFish", al cui interno sono state create le opportune repository in base alle necessità del progetto.
Le varie repository presenti nell'organizzazione sono raggiungibili al seguente link : https://github.com/orgs/SWEllfish14/repositories.

\subsection{Comandi utili di Git}
I comandi più frequentemente utilizzati sono riassunti in seguito:
\begin{itemize}
    \item Sincronizzazione con repository remota: git pull;
    \item Creazione di un nuovo branch: git branch nome branch;
    \item Passaggio a un branch specificata: git checkout nome branch;
    \item Aggiunta delle modifiche alla stage area: git add *;
    \item Creazione del commit con le modifiche: git commit;
    \item Push sul remote di un nuovo branch: git push --set-upstream origin nome branch;
    \item Push su un branch già esistente: git push
\end{itemize}

\end{document}
